\newcommand{\nomeEstado}{BAHIA}
\newcommand{\siglaEstado}{BA}
\newcommand{\nomeConcessionaria}{NEOENERGIA BAHIA}
\newcommand{\cnpjConcessionaria}{15.139.629/0001-94}
\newcommand{\termoCIP}{Sexto}
\newcommand{\presidenteConcessionaria}{Sr. Thiago Guth.}

\newcommand{\empresaResponsavel}{ABEL}
\newcommand{\nomeMunicipio}{Itabela}
\newcommand{\nomeMunicipioMaisculo}{\MakeUppercase{\nomeMunicipio}}
\newcommand{\IPestimada}{2275843}
\newcommand{\cnpjMunicipio}{16.234.429/0001-83}
\newcommand{\sedePrefeitura}{Av. Manoel Carneiro, nº 327, Centro, Itabela - BA, CEP 45848-000}
\newcommand{\nomePrefeito}{Ricardo de Jesus Flauzino}
\newcommand{\generoPrefeitura}{pelo Prefeito}
\newcommand{\numeroContrato}{018/2025}
\newcommand{\quantidadeAditivos}{0}
\newcommand{\legitimidade}{C}
\newcommand{\publicacao}{Diário Oficial do Município}
\newcommand{\dataPublicacao}{18 de agosto de 2025}
\newcommand{\tituloClausula}{CLÁUSULA PRIMEIRA: OBJETO DO CONTRATO}
\newcommand{\clausula}{CLÁUSULA PRIMEIRA: OBJETO DO CONTRATO

O CONTRATADO se obriga à contratação de pessoa jurídica para prestação de serviços técnicos especializados, visando assessorar o Município na gestão, elaboração de auditorias e laudos técnicos, mediante a conferência das faturas de energia elétrica da administração direta e indireta do Município, elaboração de memorial de cálculo de consumo e potência do parque de iluminação pública, apuração do modelo tarifário aplicado em cada unidade consumidora, bem como verificação de possíveis isenções indevidas e/ou do não repasse da Contribuição de Iluminação Pública (CIP) e/ou do não recolhimento do ISS dos prestadores de serviços do setor elétrico, visando à repetição de indébitos decorrentes de cobranças indevidas (a maior) nas contas de energia elétrica de titularidade do Município de Itabela/BA. Deverá a empresa efetuar levantamento dos créditos a que faz jus o Município, referentes aos pagamentos indevidos à concessionária de energia elétrica, conforme os prazos estabelecidos na Resolução Normativa ANEEL nº 1.000, de 7 de dezembro de 2021, especialmente o art. 323, § 2º, inciso I, e suas alterações posteriores, podendo, para tanto, representar administrativamente o Município de Itabela/BA perante a Agência Nacional de Energia Elétrica (ANEEL), perante a distribuidora de energia elétrica e, caso necessário, perante a agência reguladora conveniada do Estado, em todos os processos, reclamações, recursos e demais atos decorrentes de suas competências.}
\newcommand{\telefone}{(73) 3270-0000}
\newcommand{\email}{ouvidoria@itabela.ba.gov.br}
